\chapter{Pendahuluan}

\section{Latar Belakang}
Fenomena korelasi ekstrem antar aset seringkali mendistorsi stabilitas pasar keuangan modern secara signifikan selama periode krisis ekonomi global yang tidak terduga. Meskipun model \textit{Mean-Variance} (MV) dari \cite{markowitz_portfolio_1952} menjadi fondasi utama \textit{Modern Portfolio Theory}, asumsi distribusi \textit{return} yang normal dan stasioner membatasi reliabilitas model tersebut \cite{abdul_hali_markowitz_2020}. Ketidakmampuan model MV dalam menangkap anomali pasar serta ketergantungan non-linear menyebabkan estimasi risiko menjadi tidak akurat dan sulit diprediksi melalui komputasi klasik \cite{sikalo_efficient_2022}. Oleh karena itu, penelitian ini mentransformasikan fungsi objektif portofolio ke dalam kerangka \textit{Quadratic Unconstrained Binary Optimization} (QUBO) agar sistem komputasi kuantum dapat memproses kompleksitas finansial tersebut \cite{wang_achieving_2025}.

Para peneliti ekonofisika mengadopsi model \textit{Ising} untuk memetakan fungsi biaya portofolio ke dalam struktur Hamiltonian guna mendukung implementasi pada perangkat keras kuantum \cite{lucas_ising_2014}. Pendekatan ini merepresentasikan perilaku kolektif sistem multi-agen melalui interaksi antar variabel diskrit yang menggambarkan dinamika pasar secara lebih representatif \cite{mantegna_introduction_2000}. Representasi agen ekonomi sebagai \textit{spins} yang berinteraksi memungkinkan sistem untuk mencerminkan fenomena \textit{herding} atau sinkronisasi keputusan investor melalui keselarasan orientasi \textit{spin} \cite{hidalgo_quantum_2006}. Kerangka mekanika statistik mengabstraksikan kompleksitas pasar menjadi sistem energi, sehingga pengamatan konfigurasi energi memfasilitasi deteksi transisi fase dari kondisi stabil menuju volatilitas ekstrem.

Sistem mengonversi setiap keputusan alokasi aset portofolio menjadi variabel biner yang merepresentasikan keadaan fisik dalam model komputasi kuantum yang dibangun. Parameter medan lokal ($h_i$) menyerap informasi terkait ekspektasi imbal hasil dan preferensi risiko individual, sementara koefisien kopling ($J_{ij}$) menangkap dinamika korelasi antar aset \cite{sikalo_efficient_2022}. Penyusunan interaksi tersebut ke dalam fungsi Hamiltonian \textit{Ising} mentransformasi masalah optimasi portofolio menjadi pencarian konfigurasi energi terendah atau \textit{ground state} secara sistematis \cite{wei_solving_2025}. Logika ini memastikan bahwa solusi optimal yang dihasilkan mencerminkan titik keseimbangan energi minimum yang setara dengan efisiensi alokasi modal maksimum.

Kerangka \textit{Exact Potential Game} (EPG) mengidentifikasi interaksi strategis antar partisipan pasar dengan menawarkan pendekatan yang lebih dinamis dibandingkan ekstrapolasi statistik pada analisis deret waktu. Integrasi EPG ke dalam model \textit{Ising} memodifikasi parameter medan lokal ($h_i$) melalui perumusan fungsi potensial $\Phi$ yang diturunkan dari utilitas marginal partisipan \cite{la_torre_multi-agent_2026}. Melalui pendekatan teori permainan ini, sistem mendeteksi bias aset secara presisi dengan memodelkan setiap instrumen investasi sebagai pemain yang berinteraksi dalam pencarian titik \textit{Nash Equilibrium}. Pemanfaatan fungsi potensial menjamin bahwa setiap peningkatan utilitas individual berkontribusi langsung pada penurunan energi Hamiltonian total, sehingga menciptakan konsistensi matematis yang kokoh.

Strategi \textit{warm-start} mengatasi fenomena \textit{Barren Plateaus} yang seringkali menghambat efektivitas pembaruan parameter sirkuit variasional pada ruang pencarian yang sangat luas. Penelitian ini memanfaatkan konfigurasi biner dari \textit{Nash Equilibrium} yang dihasilkan oleh algoritma \textit{Sequential Best Response} (SBR) sebagai titik awal inisialisasi parameter VQE. Pendekatan hibrida ini memposisikan sirkuit kuantum pada wilayah probabilitas tinggi sejak iterasi pertama, sehingga mempercepat proses minimisasi energi menuju \textit{ground state}. Integrasi logika teori permainan ke dalam algoritma kuantum menjembatani dinamika strategis agen ekonomi dengan presisi komputasi kuantum dalam mengidentifikasi konfigurasi portofolio yang paling stabil.

Implementasi algoritma \textit{Variational Quantum Eigensolver} (VQE) mengadaptasi prinsip simulasi sistem kuantum untuk menyelesaikan masalah optimasi diskrit pada sektor keuangan modern. Awalnya, VQE mengestimasi \textit{ground state} pada sistem molekuler melalui minimisasi nilai ekspektasi Hamiltonian kimia yang bersifat \textit{non-commuting} secara iteratif \cite{fedorov_vqe_2022}. Transformasi fungsi objektif portofolio ke dalam model \textit{Ising} memampukan VQE untuk memetakan rezim pasar optimal melalui pencarian konfigurasi energi terendah dalam ruang QUBO \cite{wang_achieving_2025}. Algoritma ini menawarkan fleksibilitas tinggi dalam menangani keterbatasan perangkat \textit{Noisy Intermediate-Scale Quantum} (NISQ) dibandingkan dengan algoritma kuantum murni lainnya.

Kerangka kerja VQE umumnya menggunakan metode optimasi iteratif berbasis \textit{Gradient Descent} (GD) untuk memperbarui parameter menuju arah negatif gradien fungsi biaya. Meskipun GD menjadi standar dalam pembelajaran mesin klasik, implementasi pada perangkat keras kuantum menghadapi kendala \textit{sampling overhead} yang sangat besar \cite{dao_exploring_2025}. Sebagai alternatif, algoritma \textit{Simultaneous Perturbation Stochastic Approximation} (SPSA) melakukan estimasi gradien secara efisien hanya dengan dua pengukuran per iterasi terlepas dari dimensi parameter \cite{wang_achieving_2025}. SPSA menunjukkan ketangguhan yang lebih tinggi terhadap fluktuasi statistik dan \textit{noise} pengukuran yang seringkali menyebabkan ketidakstabilan pada algoritma berbasis gradien konvensional \cite{spall_implementation_1998}.

Pemilihan \textit{ansatz} yang optimal menentukan keberhasilan pencarian solusi dalam algoritma VQE secara keseluruhan pada sistem kuantum yang terbatas. \textit{Ansatz EfficientSU(2)} mengoptimalkan struktur gerbang $SU(2)$—seperti rotasi $RY$ dan $RZ$—serta gerbang \textit{CNOT} dalam pola yang sangat ringkas dan efisien \cite{dao_exploring_2025}. Dibandingkan dengan \textit{TwoLocal} yang generik, \textit{EfficientSU(2)} mereduksi kedalaman sirkuit secara signifikan tanpa mengorbankan tingkat ekspresivitas yang dibutuhkan untuk memetakan dinamika portofolio. Penggunaan \textit{ansatz} yang ringkas meminimalkan akumulasi \textit{gate noise} pada perangkat NISQ, sehingga meningkatkan akurasi estimasi energi pada setiap iterasi optimasi.

\section{Rumusan Masalah}
Berdasarkan latar belakang tersebut, permasalahan yang akan dibahas dalam tugas akhir ini adalah:
\begin{enumerate}
    \item Bagaimana perumusan bias aset sebagai parameter medan lokal ($h_i$) dan interaksi antar aset sebagai koefisien kopling ($J_{ij}$) menggunakan fungsi potensial dalam kerangka \textit{Exact Potential Game} (EPG) memengaruhi akurasi pemetaan Hamiltonian?
    \item Bagaimana efektivitas penerapan \textit{Nash Equilibrium} sebagai strategi \textit{warm-start} dalam mempercepat konvergensi dan mengatasi kendala inisialisasi pada algoritma \textit{Variational Quantum Eigensolver} (VQE)?
    \item Bagaimana performa algoritma VQE dengan optimasi \textit{Simultaneous Perturbation Stochastic Approximation} (SPSA) dan \textit{Ansatz EfficientSU(2)} dalam menemukan solusi \textit{ground state} pada model portofolio hibrida?
\end{enumerate}

\section{Tujuan Penelitian}
Dari rumusan masalah yang diberikan, tujuan dari pembuatan tugas akhir ini antara lain:
\begin{enumerate}
    \item Menganalisis formulasi bias aset sebagai parameter medan lokal ($h_i$) dan interaksi antar aset sebagai koefisien kopling ($J_{ij}$) melalui pendekatan fungsi potensial dalam sistem Hamiltonian berbasis \textit{Exact Potential Game} (EPG).
    \item Mengevaluasi kontribusi strategi \textit{warm-start} berbasis \textit{Nash Equilibrium} dalam meningkatkan efisiensi pencarian solusi optimal pada algoritma VQE.
    \item Menguji kapabilitas algoritma VQE yang menggunakan optimasi SPSA dan \textit{Ansatz EfficientSU(2)} untuk menghasilkan keputusan alokasi aset yang optimal dalam kondisi pasar yang kompleks.
\end{enumerate}

\section{Batasan Masalah}
Agar penelitian ini lebih terfokus, maka pengerjaan tugas akhir ini diberikan batasan-batasan sebagai berikut:
\begin{enumerate}
    \item Portofolio yang dianalisis dibatasi pada pemilihan 1 aset dari 2 pilihan aset dengan sektor yang berbeda.
    \item Model permainan yang digunakan adalah skema \textit{non zero-sum game} dalam kerangka \textit{Exact Potential Game}.
    \item Pencarian ekuilibrium dilakukan menggunakan algoritma \textit{Sequential Best Response} (SBR).
    \item Simulasi dan implementasi algoritma dilakukan menggunakan simulator \textit{Pennylane}.
\end{enumerate}

\section{Manfaat Penelitian}
Hasil dari tugas akhir ini diharapkan dapat memberikan alternatif metode optimasi portofolio yang lebih tangguh terhadap dinamika korelasi pasar serta memberikan kontribusi akademis dalam penerapan komputasi kuantum pada bidang finansial.
